\subsection{Wall-clock and CPU costs}
\label{sec:latency}

A directed transmission opportunity is a service budget, not a unit of
time. We therefore measure execution cost separately using the native
source revision \texttt{358b921a} from the retry-readiness experiment.
The host is a virtualized x86-64 Linux environment reporting an Intel Xeon
Platinum 8573C processor. The process is pinned to one logical CPU. We use
Rust 1.98.1 in release mode, with the default portable-C ML-KEM backend
from mlkem-native 2.0.0 and fips204 0.4.6 for ML-DSA.

Table~\ref{tab:primitive-latency} reports mean wall-clock and process CPU
time per backend call. Each operation has 256 warm-up calls and 1,600
measured calls over 64 input cases, recorded in 200 batches of eight.
Seeds, coins and input buffers are prepared before each timed interval;
signing and verification use 1,024-byte messages and an empty context.
The intervals include the backend API's key decoding and output handling,
but exclude operating-system random-number acquisition, journal work and
message transport. We retain every batch, including scheduling delays.
The artifact records the source and executable hashes, build settings,
timer control, raw wall-clock and CPU readings, and the corresponding
dispersion (Appendix~\ref{app:latency}).

% Means recomputed from the retained 358b921a raw timing batches.
\begin{table}[t]
\centering
\small
\caption{Warmed backend operation costs on the stated virtualized host.
Each mean covers 1,600 calls; all times are in milliseconds.}
\label{tab:primitive-latency}
\begin{tabularx}{\linewidth}{@{}Xrr@{}}
\toprule
Operation & Mean wall time & Mean CPU time\\
\midrule
ML-KEM-768 key generation & 0.041 & 0.040\\
ML-KEM-768 encapsulation & 0.047 & 0.047\\
ML-KEM-768 decapsulation & 0.056 & 0.056\\
ML-DSA-65 key generation & 0.381 & 0.380\\
ML-DSA-65 signing & 0.890 & 0.886\\
ML-DSA-65 verification & 0.232 & 0.232\\
\bottomrule
\end{tabularx}
\end{table}


We also measure the complete local send-to-delivery path for 256
applications across 16 independently established session pairs. Each pair
uses two separate native \texttt{FileBackend} journals, performs two warm-up
deliveries and then records 16 deliveries, alternating direction. Payloads
are 1,024 bytes with 32 bytes of associated data. The send budget is 1,024,
so this short run remains within one established epoch.
The interval starts before message-identifier allocation and ends when the
receiver returns committed plaintext. It includes the sender's reservation,
chain and outbox persistence and the receiver's authenticated processing,
chain and inbox persistence. The original protocol's random-number calls
remain inside this interval. Wire bytes pass directly between the two
endpoints in one process. Provisioning, bootstrap, application consumption
and the acknowledgement exchange occur outside it; no network delay or
rekey is included.

Table~\ref{tab:delivery-latency} reports mean wall time of 24.531\,ms and
process CPU time of 18.691\,ms. The sender and receiver portions average
16.296 and 8.236\,ms of wall time, respectively. All 256 samples are
retained, including the largest observed wall time of 257.105\,ms. These
values describe the stated host and native storage path. They do not
measure mobile scheduling, SQLite/WAL performance or the time spent
waiting for communication opportunities in the earlier workloads.

% Statistics recomputed from all 256 retained individual delivery intervals.
\begin{table}[t]
\centering
\small
\caption{Same-host delivery costs between two native file journals in an
established session. All times are in milliseconds. The 95th percentile
uses the nearest rank among all 256 individual message intervals.}
\label{tab:delivery-latency}
\begin{tabularx}{\linewidth}{@{}Xrrrr@{}}
\toprule
Interval & Mean & Median & 95th percentile & Maximum\\
\midrule
Wall-clock & 24.531 & 17.007 & 37.046 & 257.105\\
Process CPU & 18.691 & 13.591 & 27.253 & 178.035\\
\bottomrule
\end{tabularx}
\end{table}

